\documentclass[aspectratio=169,xcolor=dvipsnames]{beamer}
\usetheme{SimpleDarkBlue}
\usepackage{hyperref}
\usepackage{graphicx} % Allows including images
\usepackage{booktabs} % Allows the use of \toprule, \midrule and \bottomrule in tables
\usepackage[spanish]{babel}
\usepackage[utf8]{inputenc}
\usepackage{array}
\usepackage{tabularx}
\usepackage{caption}
\usepackage{adjustbox}
\usepackage{amssymb}
\usepackage{tcolorbox}
\usepackage{amsmath}

\graphicspath{{./graficas}} 

\title{Estadística aplicada}
\subtitle{Conceptos previos: conjuntos, funciones, cálculo y probabilidad}

\author{Práxedis Jiménez Ruvalcaba}

\institute
{

}
\date{\today}

\tcbuselibrary{theorems}

% Custom colored theorem box
\newtcbtheorem{mytheo}{Teorema}%
{colback=orange!5!white,colframe=orange!75!black,fonttitle=\bfseries}{th}

\begin{document}

\begin{frame}[c]
    \begin{beamercolorbox}[sep=8pt,center,rounded=true,shadow=true]{title}
        \usebeamerfont{title}\inserttitle\par
        \ifx\insertsubtitle\@empty
        \else
          \vskip0.25em
          {\usebeamerfont{subtitle}\insertsubtitle\par}
        \fi 
    \end{beamercolorbox}

    \vspace{1cm}
    \begin{center}
        \usebeamerfont{author}\insertauthor\par
        \vspace{0.5cm}
        \usebeamerfont{institute}\insertinstitute\par
        \usebeamerfont{date}\insertdate
    \end{center}
    
\end{frame}

\begin{section}{Cultura general de programación}

\begin{frame}{¿Por qué R?}
    \begin{enumerate}
        \item Código abierto
        \item Descendiente directo de S 
        \item RStudio
        \item Ventajas del alto lenguaje
        \item R como entorno y lenguaje de expresión
    \end{enumerate} 
\end{frame}

\begin{frame}{¿Por qué no otro?}
    \begin{itemize}
        \item JASP
        \item SAS
        \item SPSS
        \item Python
    \end{itemize}
    También existen Minitab, S-Plus, SAS, SPSS, Stata, Systat. Todos estos tienen conexión binaria con R. 
    (No es recomendado).
\end{frame}

\end{section}


\begin{section}{Notas importantes de uso}  
    
\begin{frame}{Notas sobre R}
    \begin{itemize}
        \item R nació en 1992, la documentación previa es propiamente de S (sigue siendo útil)
        \item La mayoría de documentos sobre el uso y manejo a nivel sistema están pensados para UNIX, pero hay 'traducciones' a Windows y Mac
        \item Se va a intentar cubrir todo lo posible con R base, pero se recomiendan y usarán algunas librerías CRAN
    \end{itemize}
\end{frame}

\begin{frame}{Buenas prácticas}
    \begin{itemize}
        \item Mantener estructuras de carpetas simple y consistente
        \item Siempre trabajar dentro de R Projects, separados propiamente
        \item Nombrar correctamente las variables
        \item Comentar y documentar lo mayor posible. Especialmente en limpieza de datos y elección de métodos
        \item Usar las librerías explícitamente y evitar importacion directa.
        \item Tener scripts en orden apropiado de corrida.
        \item Evitar copiar y pegar códgio
    \end{itemize}
\end{frame}

\begin{frame}{Buenas Prácticas: Evitar copiar y pegar}
    Al presentarse dudas o errores, ¿qué debería hacerse?
    \begin{enumerate}
        \item Leer documentación o autodocumentación
        \item 'probar' y 'moverle'
        \item Buscar en foros
        \item Preguntar a alguien
        \item Preguntar a la IA
    \end{enumerate}
\end{frame}
\end{section}

\begin{section}{Notas de uso propias de R}
\begin{frame}{Algunas convenciones de lenguaje}
    Los datos "nulos" pueden presentarse de varias formas según el tipo de falta
    \begin{itemize}
        \item NA   $->$ datos faltantes
        \item NaN  $->$ not a number
        \item NULL $->$ Ausencia de valor
    \end{itemize}
\end{frame}

\begin{frame}{Algunas convenciones del lenguaje}
    \begin{itemize}
        \item Procurar usar la asignación "$<-$"
        \item Para escribir strings/char usar dobles comillas ''char''
        \item La indexación comienza en 1
        \item Los logaritmos si no se especifica, son neperianos
        \item Tener cuidado con la notación en la documentación
    \end{itemize}  
\end{frame}

\begin{frame}{Salvavidas}
    \begin{itemize}
        \item Recordar al inicio y al final de todo proyecto, usar el comando \textit{SessionInfo()}
        \item El mayor salvavidas para acceder a documentación de cualquien función \textit{function} es \textit{?function} y \textit{??function}
        \item .RDAta y .Rhistory 
    \end{itemize}

\end{frame}

\end{section}



\end{document}